\documentclass[11pt]{article}

\usepackage[a4paper,margin=2cm]{geometry}
\usepackage[utf8]{inputenc}
\usepackage[T1]{fontenc}
\usepackage{amsmath}
\usepackage{amssymb}

\newcommand{\IR}{\mathbb{R}}
\newcommand{\norm}[1]{\left\lVert #1 \right\rVert}

\title{Why a Steepest-Descent Line Search at $k=0$ is Consistent with\\
\emph{Minimization of expensive functions in the presence of forbidden
regions}}
\author{}
\date{}

\begin{document}
\maketitle

\section{The degeneracy at $k=0$}

\texttt{ffjm2} minimizes $f(x) = \tfrac12\norm{F(x)}^2$, where
$F=(f_1,\dots,f_m)^\top$ is the residual vector. At external iteration $k$
it keeps one symmetric matrix $H_i^k$ per residual component $f_i$,
approximating $\nabla^2 f_i(x^k)$, and solves
\begin{equation}
    M_k(s) \;=\; \sigma\norm{s}^2 + \tfrac12\sum_i q_i(s)^2, \qquad
    q_i(s) = f_i(x^k) + \nabla f_i(x^k)^\top s + \tfrac12 s^\top H_i^k s,
    \label{eq:quartic}
\end{equation}
accepting or rejecting $x^k+s$ by the ratio between actual and predicted
reduction, and growing/shrinking $\sigma$ accordingly (this is the
regularization scheme of Algorithm 4.1 of \cite{FortunatoMartinez2026},
implemented in \texttt{ffjm2} with an empirical persistent-$\sigma$
variant --- the paper resets $\sigma\leftarrow0$ at the start of every
iteration and never shrinks it, but persisting it between iterations and
shrinking it on acceptance performed better in practice).

At $k=0$, every $H_i^0=0$ (no secant pair has been observed yet), so
\eqref{eq:quartic} collapses to
\begin{equation}
    M_0(s) = \sigma\norm{s}^2 + \tfrac12\sum_i \bigl(f_i(x^0) + \nabla f_i(x^0)^\top s\bigr)^2,
    \label{eq:gn}
\end{equation}
a single damped Gauss--Newton subproblem. Two things are worth noting
about this particular instance of the subproblem:
\begin{enumerate}
    \item Unlike every subsequent iteration, $\sigma=\sigma_0$ here is a
    fixed constant, not yet informed by any accept/reject test on
    \emph{this} problem. Whether it makes \eqref{eq:gn} behave like
    Gauss--Newton ($\sigma\to0$) or like scaled steepest descent
    ($\sigma\to\infty$, where the minimizer of \eqref{eq:gn} is
    asymptotically parallel to $-\nabla f(x^0)$) is essentially a guess
    about the local curvature scale of $f$ at $x^0$ that the method has
    had no chance to calibrate yet.
    \item There is, by construction, no curvature information at all to
    justify solving a quartic (rather than a quadratic) model here: the
    quartic term $\tfrac12 s^\top H_i^k s$ is identically zero for every
    $i$.
\end{enumerate}
So the very first step is the one iteration where the method's own
machinery (quasi-Newton curvature, calibrated regularization) has nothing
to work with.

\section{A self-contained argument: the closed form of $M_0$}
\label{sec:closedform}

Rather than appealing to a different algorithm, the cleanest justification
comes directly out of \eqref{eq:gn} itself, which is an unconstrained,
strictly convex quadratic in $s$ and can be solved in closed form. Write
$J:=F'(x^0)$ for the Jacobian of $F$ at $x^0$ (the matrix whose $i$-th row
is $\nabla f_i(x^0)^\top$), so that $g:=\nabla f(x^0)=J^\top F(x^0)$, and
$A:=J^\top J$ (symmetric positive semidefinite). Then
\[
    \nabla_s M_0(s) = (2\sigma I + A)\,s + g = 0
    \quad\Longrightarrow\quad
    s(\sigma) = -(A+2\sigma I)^{-1} g.
\]

\textbf{Lemma A (descent for every $\sigma>0$).} Since $A+2\sigma I \succ
0$ for $\sigma>0$, its inverse is symmetric positive definite, hence
\[
    g^\top s(\sigma) = -g^\top (A+2\sigma I)^{-1} g < 0 \qquad \text{whenever } g\ne0.
\]
So $s(\sigma)$ is a genuine descent direction for $f$ \emph{for every
$\sigma>0$}, not just asymptotically --- the $\sigma$-loop never has a
``wrong-direction'' problem at $k=0$; what it can have is a
badly-\emph{sized} step, which is what the accept/reject test and the
growth of $\sigma$ are actually correcting for.

\textbf{Proposition B (alignment with $-g$ as $\sigma\to\infty$).} Factor
$2\sigma$ out of the inverse:
\[
    s(\sigma) = -(A+2\sigma I)^{-1}g
    = -\frac{1}{2\sigma}\left(I + \frac{A}{2\sigma}\right)^{-1}g.
\]
As $\sigma\to\infty$, $A/(2\sigma)\to0$, and matrix inversion is
continuous at the identity, so $\bigl(I+A/(2\sigma)\bigr)^{-1}\to I$.
Hence
\begin{equation}
    s(\sigma) = -\frac{1}{2\sigma}\,g + o\!\left(\frac1\sigma\right)
    \qquad\Longrightarrow\qquad
    \frac{s(\sigma)}{\norm{s(\sigma)}} \longrightarrow -\frac{g}{\norm{g}}
    \quad\text{as } \sigma\to\infty.
    \label{eq:asymptotic}
\end{equation}

In words: \emph{for $\sigma$ large enough, the solution of the degenerate
$k=0$ subproblem \eqref{eq:gn} already is, to leading order, a
steepest-descent step of length $\approx1/2\sigma$} --- the two are
asymptotically the same computation. This is a direct consequence of
\eqref{eq:gn}'s own closed form, entirely internal to Algorithm 4.1.

In practice, the rejection ladder $\sigma_0,\ \sigma_{j+1}=\gamma\sigma_j$
($\gamma>1$ the growth factor) tries $s(\sigma_0), s(\sigma_1),\dots$
until one is accepted. By Proposition B, once $\sigma_j\gg\norm{A}/2$ the
successive trials are already essentially parallel to $-g$ and shrink by
the fixed factor $1/\gamma$ --- so once it reaches that scale, the ladder
is already behaving like a geometric backtracking line search along $-g$.
Below that scale, though --- which includes $\sigma_0$ itself whenever
$\sigma_0$ (a fixed constant, unrelated to the problem at hand) is not
already $\gg\norm{A}/2$ --- $s(\sigma_0)$ is a perfectly valid descent
direction (Lemma A) but generally \emph{not} aligned with $-g$, and its
length is not calibrated to $\norm{g}$ at all: whether it is accepted on
the first try is essentially a coincidence of scale, not a controlled
property of the method.

Replacing the $k=0$ subproblem by an explicit Armijo line search along
$-g$ therefore does not change what the method is doing once $\sigma$ is
large --- it is, to leading order, the same computation. What it removes
is the uncalibrated, possibly wasted, low-$\sigma$ rejections at the start
of the ladder, by searching directly along the direction
\eqref{eq:asymptotic} shows the ladder is heading toward anyway, with a
step-length search (quadratic interpolation) that is more
sample-efficient than a fixed geometric factor $1/\gamma$. It is also
computationally cheaper: each rejected $\sigma$ in the ladder re-invokes
the general-purpose numerical solver used for the full quartic subproblem
(which explores several random restarts to guard against spurious local
minima), even though at $H_i=0$ the subproblem has the trivial closed
form $s(\sigma)=-(A+2\sigma I)^{-1}g$ above --- a backtracking line search
only needs cheap residual evaluations $F(x^0+\alpha p)$ to test
sufficient decrease.


\begin{thebibliography}{9}

\bibitem{FortunatoMartinez2026}
F. A. Fortunato and J. M. Mart\'inez, \emph{Minimization of expensive
functions in the presence of forbidden regions}, unpublished manuscript,
2026. (\texttt{\_research/ffjm2.pdf} in this repository.)

\end{thebibliography}

\end{document}
